\chapter{Random Vector}

\section{Bivariate Random Variables}

\begin{definition}[Bivariate Random Variable]\index{bivariate random variable}
  A pair of random variables \((X, Y)\) defined on the same sample space \(\Omega\) is called a \emph{bivariate random variable}.
\end{definition}

\begin{definition}[Joint Distribution Function]\index{joint distribution function}
  The \emph{joint distribution function} (joint d.f.\@) of a bivariate random variable \((X, Y)\) is defined as
  \[
    F(x, y) = P\!\left(\{X \leq x\} \cap \{Y \leq y\}\right) = P(X \leq x, Y \leq y).
  \]
\end{definition}

\begin{proposition}
  Some properties of the joint d.f.\@ \(F(x, y)\) of a bivariate random variable \((X, Y)\) are as follows:
  \begin{enumerate}
    \item (Monotonicity) If \(y_2 > y_1\), then \(F(x, y_2) \geq F(x, y_1)\) for any \(x\); if \(x_2 > x_1\), then \(F(x_2, y) \geq F(x_1, y)\) for any \(y\).
    \item (Right-continuity) \(\lim_{x \to x_0^+} F(x, y) = F(x_0, y)\) for any \(y\); \(\lim_{y \to y_0^+} F(x, y) = F(x, y_0)\) for any \(x\).
    \item \(0 \leq F(x, y) \leq 1\) for all \(x, y \in \mathbb{R}\), and
    \begin{align*}
      F(x, -\infty) = \lim_{y \to -\infty} F(x, y) = 0, &\quad F(-\infty, y) = \lim_{x \to -\infty} F(x, y) = 0, \\
      F(-\infty, -\infty) = \lim_{x \to -\infty} \lim_{y \to -\infty} F(x, y) = 0, &\quad F(\infty, \infty) = \lim_{x \to \infty} \lim_{y \to \infty} F(x, y) = 1.
    \end{align*}
    \item (Non-negativity) For any \(x_1 < x_2\) and \(y_1 < y_2\), we have
    \begin{align*}
      &P(x_1 < X \leq x_2, y_1 < Y \leq y_2) \\
      &= F(x_2, y_2) - F(x_2, y_1) - F(x_1, y_2) + F(x_1, y_1) \geq 0.
    \end{align*}
  \end{enumerate}
\end{proposition}

\begin{example}
  The joint d.f.\@ of a bivariate random variable \((X, Y)\) is given by
  \[
    F(x, y) = A \left(B + \arctan\frac{x}{2}\right) \left(C + \arctan\frac{y}{2}\right), \quad x, y \in \mathbb{R}.
  \]
  \begin{enumerate}
    \item Find the values of the constants \(A, B, C\).
    \item Determine the value of \(P(0 < X \leq 2, 0 < Y \leq 3)\).
  \end{enumerate}
\end{example}
\begin{solution}
  \begin{enumerate}
    \item Since \(F(x, y)\) is a joint d.f.\@, it must satisfy the properties stated in the previous proposition. In particular, we have
    \[
      \begin{cases}
        F(-\infty, y) = 0 \implies A (B - \frac{\pi}{2}) (C + \arctan\frac{y}{2}) = 0, \\
        F(x, -\infty) = 0 \implies A (B + \arctan\frac{x}{2}) (C - \frac{\pi}{2}) = 0, \\
        F(\infty, \infty) = 1 \implies A (B + \frac{\pi}{2}) (C + \frac{\pi}{2}) = 1.
      \end{cases}
    \]
    Solving these equations, we find \(A = \frac{1}{\pi^2}\), \(B = \frac{\pi}{2}\), and \(C = \frac{\pi}{2}\).
    \item We have
    \[
      P(0 < X \leq 2, 0 < Y \leq 3) = F(2, 3) - F(2, 0) - F(0, 3) + F(0, 0)
    \]
    Substituting the values of \(A, B, C\) and evaluating the joint d.f.\@ at the specified points, we can compute the desired probability. \qedhere
  \end{enumerate}
\end{solution}

\begin{definition}[Joint Probability Function]\index{joint probability function}
  For discrete bivariate random vectors \((X, Y)\), whose possible values are a countable set of points \((x_i, y_j) (i, j = 1, 2, \dots)\), the \emph{joint probability function} (joint p.f.\@) is defined as
  \[
    p(x, y) = P(X = x, Y = y).
  \]
\end{definition}

\begin{proposition}
  For a discrete bivariate random variable \((X, Y)\) with joint p.f.\@ \(p_{ij} = P(X = x_i, Y = y_j)\), we have
  \begin{enumerate}
    \item \(p_{ij} \geq 0\) for all \(i, j\).
    \item \(\sum_{i=1}^\infty \sum_{j=1}^\infty p_{ij} = 1\).
    \item The joint d.f.\@ \(F(x, y)\) can be expressed in terms of the joint p.f.\@ as
    \[
      F(x, y) = P(X \leq x, Y \leq y) = \sum_{x_i \leq x} \sum_{y_j \leq y} p_{ij}.
    \]
  \end{enumerate}
\end{proposition}

\begin{example}
  Assume that \(Y \sim \mathrm{N}(0, 1)\). Find the joint p.f.\@ of
  \[
    X_1 = \begin{cases}
      0, & |Y| \geq 1, \\
      1, & |Y| < 1,
    \end{cases}
    X_2 = \begin{cases}
      0, & |Y| \geq 2, \\
      1, & |Y| < 2.
    \end{cases}
  \]
\end{example}
\begin{solution}
  \begin{align*}
    P(X_1 = 0, X_2 = 0) &= P(|Y| \geq 2) = 2P(Y \geq 2) = 2(1 - \Phi(2)), \\
    P(X_1 = 0, X_2 = 1) &= P(|Y| \geq 1, |Y| < 2) = P(1 \leq |Y| < 2) = 2(\Phi(2) - \Phi(1)), \\
    P(X_1 = 1, X_2 = 0) &= P(|Y| < 1, |Y| \geq 2) = 0, \\
    P(X_1 = 1, X_2 = 1) &= P(|Y| < 1) = 2\Phi(1) - 1. \qedhere
  \end{align*}
\end{solution}

\begin{definition}[Joint Probability Density Function]\index{joint probability density function}\index{jointly continuous}
  For a bivariate random variable \((X, Y)\), if there exists a non-negative function \(f(x, y)\) such that
  \[
    F(x, y) = P(X \leq x, Y \leq y) = \int_{-\infty}^x \int_{-\infty}^y f(u, v) \, \mathrm{d}u \, \mathrm{d}v,
  \]
  for all \(x, y \in \mathbb{R}\), then \(X\) and \(Y\) are called \emph{jointly continuous}, and the function \(f(x, y)\) is called the \emph{joint probability density function} (joint p.d.f.\@) of \((X, Y)\).
\end{definition}

\begin{proposition}
  For a jointly continuous bivariate random variable \((X, Y)\) with joint p.d.f.\@ \(f(x, y)\), we have
  \begin{enumerate}
    \item \(f(x, y) \geq 0\) for all \(x, y \in \mathbb{R}\).
    \item \(\iint_{\mathbb{R}^2} f(x, y) \, \mathrm{d}x \, \mathrm{d}y = 1\).
    \item If \(f(x, y)\) is continuous at a point \((x, y)\), then
    \[
      f(x, y) = \frac{\partial^2 F(x, y)}{\partial x \partial y} = \lim_{\Delta x,\Delta y\to0^+}\frac{P(x<X\le x+\Delta x,\ y<Y\le y+\Delta y)}{\Delta x,\Delta y}.
    \]
    \item For any \(G \subseteq \mathbb{R}^2\), we have
    \[
      P\{(X, Y) \in G\} = \iint_G f(x, y) \, \mathrm{d}x \, \mathrm{d}y.
    \]
  \end{enumerate}
\end{proposition}

\begin{definition}[Bivariate Uniform Distribution]\index{bivariate uniform distribution}
  A bivariate random variable \((X, Y)\) is said to have a \emph{bivariate uniform distribution} on a measurable region \(R\) with finite positive area if its joint p.d.f.\@ is given by
  \[
    f(x, y) = \begin{cases}
      \frac{1}{\text{area}(R)}, & (x, y) \in R, \\
      0, & (x, y) \notin R.
    \end{cases}
  \]
\end{definition}

\begin{definition}[Bivariate Normal Distribution]\index{bivariate normal distribution}
  A bivariate random variable \((X, Y)\) is said to have a \emph{bivariate normal distribution} with parameters \(\mu_1, \mu_2, \sigma_1 > 0, \sigma_2 > 0\), and \(|\rho| < 1\) if its joint p.d.f.\@ is given by
  \[
    f(x, y) = \frac{1}{2\pi \sigma_1 \sigma_2 \sqrt{1 - \rho^2}} e^{-\frac{1}{2(1 - \rho^2)} \left[\frac{(x - \mu_1)^2}{\sigma_1^2} - 2\rho \frac{(x - \mu_1)(y - \mu_2)}{\sigma_1 \sigma_2} + \frac{(y - \mu_2)^2}{\sigma_2^2}\right]}.
  \]
  We denote this distribution by \((X, Y) \sim \mathrm{N}(\mu_1, \mu_2; \sigma_1^2, \sigma_2^2; \rho)\).
\end{definition}

\begin{example}
  Suppose that \((X, Y) \sim f(x, y)\) with
  \[
    f(x, y) = \begin{cases}
      A e^{-(2x + 3y)}, & x \geq 0, y \geq 0 \\
      0, & \text{otherwise}.
    \end{cases}
  \]
  \begin{enumerate}
    \item Find the value of the constant \(A\).
    \item Calculate the probability \(P(X < 2, Y < 1)\).
  \end{enumerate}
\end{example}
\begin{solution}
  \begin{enumerate}
    \item Since \(f(x, y)\) is a joint p.d.f.\@, it must satisfy the property that \(\iint_{\mathbb{R}^2} f(x, y) \, \mathrm{d}x \, \mathrm{d}y = 1\). Therefore, we have
    \begin{align*}    
      \iint_{\mathbb{R}^2} f(x, y) \, \mathrm{d}x \, \mathrm{d}y &= \int_0^\infty \int_0^\infty A e^{-(2x + 3y)} \, \mathrm{d}x \, \mathrm{d}y \\
      &= A \int_0^\infty e^{-2x} \, \mathrm{d}x \int_0^\infty e^{-3y} \, \mathrm{d}y \\
      &= A \left(-\frac{1}{2} e^{-2x}\right)\Big|_0^\infty \left(-\frac{1}{3} e^{-3y}\right)\Big|_0^\infty = \frac{1}{6} A = 1.
    \end{align*}
    Solving for \(A\), we find \(A = 6\).
    \item We can compute the desired probability as follows:
    \begin{align*}
      P(X < 2, Y < 1) &= \iint_{0 \leq x < 2, 0 \leq y < 1} f(x, y) \, \mathrm{d}x \, \mathrm{d}y \\
      &= \int_0^2 \int_0^1 6 e^{-(2x + 3y)} \, \mathrm{d}y \, \mathrm{d}x = 6 \int_0^2 e^{-2x} \, \mathrm{d}x \int_0^1 e^{-3y} \, \mathrm{d}y \\
      &= 6 \left(-\frac{1}{2} e^{-2x}\right)\Big|_0^2 \left(-\frac{1}{3} e^{-3y}\right)\Big|_0^1 = (1 - e^{-3})(1 - e^{-4}). \qedhere
    \end{align*}
  \end{enumerate}
\end{solution}

\begin{example}
  Suppose that the joint p.d.f.\@ of \(X, Y\) is given by
  \[
    f(x, y) = \begin{cases}
      e^{-y}, & 0 < x < y \\
      0, & \text{otherwise}.
    \end{cases}
  \]
  Calculate the probability \(P(X + Y \leq 1)\).
\end{example}
\begin{solution}
  \begin{align*}
    P(X + Y \leq 1) &= \iint_{x + y \leq 1} f(x, y) \, \mathrm{d}x \, \mathrm{d}y \\
    &= \int_0^\frac{1}{2} \mathrm{d}x \int_x^{1 - x} e^{-y} \, \mathrm{d}y \\
    &= 1 + e^{-1} - 2e^{-\frac{1}{2}}. \qedhere
  \end{align*}
\end{solution}

\section{Marginal Distribution}

\begin{definition}[Marginal Distribution]\index{marginal distribution}
  Suppose that \(F(x, y)\) is the joint d.f.\@ of a bivariate random variable \((X, Y)\). The \emph{marginal distributions} of \(X\) and \(Y\) are defined as
  \[
    F_X(x) = P(X \leq x) = F(x, \infty), \quad F_Y(y) = P(Y \leq y) = F(\infty, y).
  \]
\end{definition}

In the case of discrete bivariate random variables with joint p.f.\@ \(p_{ij} = P(X = x_i, Y = y_j)\), the marginal distributions can be expressed in terms of the joint p.f.\@ as
\[
  F_X(x) = P(X \leq x) = \sum_{x_i \leq x} \sum_{j=1}^\infty p_{ij}, \quad F_Y(y) = P(Y \leq y) = \sum_{i=1}^\infty \sum_{y_j \leq y} p_{ij}.
\]
The marginal p.f.\@ of \(X\) and \(Y\) are given by
\[
  p_{i \cdot} = P(X = x_i) = \sum_{j=1}^\infty p_{ij}, \quad p_{\cdot j} = P(Y = y_j) = \sum_{i=1}^\infty p_{ij}.
\]

In the case of jointly continuous bivariate random variables, the marginal distributions can be expressed in terms of the joint p.d.f.\@ as
\begin{align*}
  F_X(x) &= P(X \leq x) = \int_{-\infty}^x \left[\int_{-\infty}^\infty f(u, y) \, \mathrm{d}y \right] \mathrm{d}u, \\
  F_Y(y) &= P(Y \leq y) = \int_{-\infty}^y \left[\int_{-\infty}^\infty f(x, v) \, \mathrm{d}x \right] \mathrm{d}v.
\end{align*}
The marginal p.d.f.\@ of \(X\) and \(Y\) is given by
\[
  f_X(x) = \int_{-\infty}^\infty f(x, y) \, \mathrm{d}y, \quad f_Y(y) = \int_{-\infty}^\infty f(x, y) \, \mathrm{d}x.
\]

\begin{example}
  Suppose that \(X, Y\) have the uniform distribution on the region \(D_1 = \{(x, y), x^2 + y^2 \leq 1\}\). Find the marginal p.d.f.\@ of \(X\).
\end{example}
\begin{solution}
  Since \(X, Y\) have the uniform distribution on the region \(D_1\), we have
  \[
    f(x, y) = \begin{cases}
      \frac{1}{\pi}, & x^2 + y^2 \leq 1, \\
      0, & \text{otherwise}.
    \end{cases}
  \]
  \begin{enumerate}
    \item When \(|x| > 1\), \(f(x, y) = 0\) for all \(y\), so \(f_X(x) = 0\).
    \item When \(|x| \leq 1\), we have
    \[
      f_X(x) = \int_{-\sqrt{1 - x^2}}^{\sqrt{1 - x^2}} \frac{1}{\pi} \, \mathrm{d}y = \frac{2}{\pi}\sqrt{1 - x^2}.
    \]
  \end{enumerate}
  That is,
  \[
    f_X(x) = \begin{cases}
      \frac{2}{\pi}\sqrt{1 - x^2}, & -1 \leq x \leq 1, \\
      0, & \text{otherwise}.
    \end{cases} \qedhere
  \]
\end{solution}

\begin{proposition}
  If \((X, Y) \sim \mathrm{U}([a, b] \times [c, d])\), then \(X \sim \mathrm{U}(a, b)\) and \(Y \sim \mathrm{U}(c, d)\).
\end{proposition}

\begin{proposition}
  If \((X, Y) \sim \mathrm{N}(\mu_1, \mu_2; \sigma_1^2, \sigma_2^2; \rho)\), then \(X \sim \mathrm{N}(\mu_1, \sigma_1^2)\) and \(Y \sim \mathrm{N}(\mu_2, \sigma_2^2)\).
\end{proposition}
\begin{proof}
  Since \((X, Y) \sim \mathrm{N}(\mu_1, \mu_2; \sigma_1^2, \sigma_2^2; \rho)\), the joint p.d.f.\@ of \(X\) and \(Y\) is given by
  \[
    f(x, y) = \frac{1}{2\pi \sigma_1 \sigma_2 \sqrt{1 - \rho^2}} e^{-\frac{1}{2(1 - \rho^2)} \left[\frac{(x - \mu_1)^2}{\sigma_1^2} - 2\rho \frac{(x - \mu_1)(y - \mu_2)}{\sigma_1 \sigma_2} + \frac{(y - \mu_2)^2}{\sigma_2^2}\right]}.
  \]
  Since
  \[
    \frac{(y - \mu_2)^2}{\sigma_2^2} - 2\rho \frac{(x - \mu_1)(y - \mu_2)}{\sigma_1 \sigma_2} = \left(\frac{y - \mu_2}{\sigma_2} - \rho \frac{x - \mu_1}{\sigma_1}\right)^2 - \rho^2 \frac{(x - \mu_1)^2}{\sigma_1^2},
  \]
  we can calculate the marginal p.d.f.\@ of \(X\) as follows:
  \begin{align*}
    f_X(x) &= \int_{-\infty}^\infty f(x, y) \, \mathrm{d}y \\
    &= \int_{-\infty}^\infty \frac{1}{2\pi \sigma_1 \sigma_2 \sqrt{1 - \rho^2}} e^{-\frac{1}{2(1 - \rho^2)} \left[\frac{(x - \mu_1)^2}{\sigma_1^2} + \left(\frac{y - \mu_2}{\sigma_2} - \rho \frac{x - \mu_1}{\sigma_1}\right)^2 - \rho^2 \frac{(x - \mu_1)^2}{\sigma_1^2}\right]} \, \mathrm{d}y \\
    &= \frac{1}{2\pi \sigma_1 \sigma_2 \sqrt{1 - \rho^2}} e^{-\frac{(x - \mu_1)^2}{2\sigma_1^2}} \int_{-\infty}^\infty e^{-\frac{1}{2(1 - \rho^2)}\left(\frac{y - \mu_2}{\sigma_2} - \rho \frac{x - \mu_1}{\sigma_1}\right)^2} \, \mathrm{d}y.
  \end{align*}
  Let \(t = \frac{1}{\sqrt{1 - \rho^2}} \left(\frac{y - \mu_2}{\sigma_2} - \rho \frac{x - \mu_1}{\sigma_1}\right)\), then \(\mathrm{d}t = \frac{1}{\sqrt{1 - \rho^2}} \frac{1}{\sigma_2} \, \mathrm{d}y\). Therefore,
  \[
    f_X(x) = \frac{1}{2\pi \sigma_1} e^{-\frac{(x - \mu_1)^2}{2\sigma_1^2}} \int_{-\infty}^\infty e^{-\frac{t^2}{2}} \, \mathrm{d}t = \frac{1}{\sqrt{2\pi} \sigma_1} e^{-\frac{(x - \mu_1)^2}{2\sigma_1^2}}.
  \]
  So \(X \sim \mathrm{N}(\mu_1, \sigma_1^2)\). By symmetry, we can also show that \(Y \sim \mathrm{N}(\mu_2, \sigma_2^2)\).
\end{proof}

\section{Conditional Distribution}

\begin{definition}[Conditional Probability Function]\index{conditional probability function}
  For a discrete bivariate random variable \((X, Y)\) with joint p.f.\@ \(p_{ij} = P(X = x_i, Y = y_j)\), the \emph{conditional probability function} (conditional p.f.\@) of \(X\) given \(Y = y_j\) is defined as
  \[
    p_{i|j} = P(X = x_i | Y = y_j) = \frac{P(X = x_i, Y = y_j)}{P(Y = y_j)} = \frac{p_{ij}}{p_{\cdot j}},
  \]
  provided that \(p_{\cdot j} > 0\).
\end{definition}

\begin{definition}[Conditional Distribution Function]\index{conditional distribution function}
  For a discrete bivariate random variable \((X, Y)\) with joint p.f.\@ \(p_{ij} = P(X = x_i, Y = y_j)\), the \emph{conditional distribution function} (conditional d.f.\@) of \(X\) given \(Y = y_j\) is defined as
  \[
    F_{X | Y}(x | y_j) = P\{X \leq x | Y = y_j\} = \frac{\sum_{x_i \leq x} p_{ij}}{p_{\cdot j}}
  \]
  provided that \(p_{\cdot j} > 0\).
\end{definition}

\begin{definition}[Conditional Probability Density Function]\index{conditional probability density function}
  For a jointly continuous bivariate random variable \((X, Y)\) with joint p.d.f.\@ \(f(x, y)\), the \emph{conditional probability density function} (conditional p.d.f.\@) of \(X\) given \(Y = y\) is defined as
  \[
    f_{X|Y}(x | y) = \frac{f(x, y)}{f_Y(y)},
  \]
  where \(f_Y(y)\) is the marginal p.d.f.\@ of \(Y\), provided that \(f_Y(y) > 0\).
\end{definition}

For a continuous random variable $Y$, the probability of exactly observing $Y = y$ is zero ($P(Y=y) = 0$). To define the conditional distribution of $X$ given $Y=y$, we must therefore consider an interval around $y$ and take the limit as the interval length goes to zero to capture the behavior of $X$ as $Y$ approaches $y$.

\begin{definition}[Conditional Distribution Function]\index{conditional distribution function}
  In the case of a jointly continuous bivariate random variable \((X, Y)\) with joint p.d.f.\@ \(f(x, y)\), for any \(x\), the limit
  \[
    F_{X | Y}(x | y) = \lim_{\varepsilon \to 0^+} P\{X \leq x | y - \varepsilon < Y \leq y + \varepsilon\} = \int_{-\infty}^{x} \frac{f(u, y)}{f_Y(y)} \, \mathrm{d}u
  \]
  is called \emph{conditional distribution function} (conditional d.f.\@) of \(X\) given \(Y = y\) if it exists, provided that \(f_Y(y) > 0\).
\end{definition}

\begin{example}
  Suppose that the number \(X\) is taken randomly from the interval \((0, 1)\). When \(X = x (0 < x < 1)\) is given, the number \(Y\) is taken randomly from the interval \((x, 1)\). Find the marginal p.d.f.\@ \(f_Y(y)\) of \(Y\).
\end{example}
\begin{solution}
  Since \(X\) is taken randomly from the interval \((0, 1)\), the p.d.f.\@ of \(X\) is given by
  \[
    f_X(x) = \begin{cases}
      1, & 0 < x < 1, \\
      0, & \text{otherwise}.
    \end{cases}
  \]
  For any \(0 < x < 1\), when \(X = x\) is given, the p.d.f.\@ of \(Y\) is given by
  \[
    f_{Y|X}(y | x) = \begin{cases}
      \frac{1}{1 - x}, & x < y < 1, \\
      0, & \text{otherwise}.
    \end{cases}
  \]
  Therefore, the joint p.d.f.\@ of \(X\) and \(Y\) is given by
  \[
    f(x, y) = f_X(x) f_{Y|X}(y | x) = \begin{cases}
      \frac{1}{1 - x}, & 0 < x < y < 1, \\
      0, & \text{otherwise}.
    \end{cases}
  \]
  Thus, the marginal p.d.f.\@ of \(Y\) is given by
  \[
    f_Y(y) = \int_{-\infty}^\infty f(x, y) \, \mathrm{d}x = \begin{cases}
      \int_0^y \frac{1}{1 - x} \, \mathrm{d}x = -\ln(1 - y), & 0 < y < 1, \\
      0, & \text{otherwise}.
    \end{cases} \qedhere
  \]
\end{solution}

\begin{example}
  Suppose that the joint p.d.f.\@ of \(X, Y\) is given by
  \[
    f(x, y) = \begin{cases}
      3x, & 0 < x < 1, 0 < y < x, \\
      0, & \text{otherwise}.
    \end{cases}
  \]
  \begin{enumerate}
    \item Find the marginal p.d.f.\@ of \(X\) and \(Y\).
    \item Find the conditional p.d.f.\@ of \(X\) and \(Y\).
  \end{enumerate}
\end{example}
\begin{solution}
  \begin{enumerate}
    \item The marginal p.d.f.\@ of \(X\) is given by
    \[
      f_X(x) = \int_{-\infty}^\infty f(x, y) \, \mathrm{d}y = \int_0^x 3x \, \mathrm{d}y = 3x^2, \quad 0 < x < 1,
    \]
    and the marginal p.d.f.\@ of \(Y\) is given by
    \[
      f_Y(y) = \int_{-\infty}^\infty f(x, y) \, \mathrm{d}x = \int_y^1 3x \, \mathrm{d}x = \frac{3}{2}(1 - y^2), \quad 0 < y < 1.
    \]
    For other values of \(x\) and \(y\), the marginal p.d.f.\@s are zero.
    \item If \(0 < y < 1\), the conditional p.d.f.\@ of \(X\) given \(Y = y\) is given by
    \[
      f_{X|Y}(x | y) = \frac{f(x, y)}{f_Y(y)} = \begin{cases}
        \frac{2x}{1 - y^2}, & y < x < 1, \\
        0, & \text{otherwise}.
      \end{cases}
    \]
    If \(0 < x < 1\), the conditional p.d.f.\@ of \(Y\) given \(X = x\) is given by
    \[
      f_{Y|X}(y | x) = \frac{f(x, y)}{f_X(x)} = \begin{cases}
        \frac{1}{x}, & 0 < y < x, \\
        0, & \text{otherwise}.
      \end{cases} \qedhere
    \]
  \end{enumerate}
\end{solution}

\begin{proposition}
  Suppose that \((X, Y) \sim \mathrm{N}(\mu_1, \mu_2; \sigma_1^2, \sigma_2^2; \rho)\). Then the conditional distribution of \(X\) given \(Y = y\) is a normal distribution with mean \(\mu_1 + \rho \frac{\sigma_1}{\sigma_2} (y - \mu_2)\) and variance \((1 - \rho^2) \sigma_1^2\).
\end{proposition}

\section{Independence of Variables}

\begin{definition}[Independence]\index{independence}
  A bivariate random variable \((X, Y)\) is said to be \emph{independent} if for any \(x, y \in \mathbb{R}\), we have
  \[
    P(X \leq x, Y \leq y) = P(X \leq x) P(Y \leq y),
  \]
  or equivalently,
  \[
    F(x, y) = F_X(x) F_Y(y).
  \]
\end{definition}

\begin{proposition}
  A discrete bivariate random variable \((X, Y)\) with joint p.f.\@ \(p_{ij} = P(X = x_i, Y = y_j)\) is independent if and only if there exist two sequences of non-negative numbers \(\{p_{i \cdot}\}\) and \(\{p_{\cdot j}\}\) such that \(p_{ij} = p_{i \cdot} p_{\cdot j}\) for all \(i, j\).
\end{proposition}
\begin{proof}
  \((\Rightarrow)\) Suppose that \((X, Y)\) is independent. Then for any \(x_1 < x_2\) and \(y_1 < y_2\), we have
  \[
    P\{x_1 < X \leq x_2, y_1 < Y \leq y_2\} = P\{x_1 < X \leq x_2\} P\{y_1 < Y \leq y_2\}.
  \]
  For any \(i, j\), we have
  \begin{align*}
    p_{ij} &= P\{X = x_i, Y = y_j\} \\
    &= \lim_{n, m \to \infty} P\{x_i - \frac{1}{n} < X \leq x_i, y_j - \frac{1}{m} < Y \leq y_j\} \\
    &= \lim_{n \to \infty} P\{x_i - \frac{1}{n} < X \leq x_i\} \lim_{m \to \infty} P\{y_j - \frac{1}{m} < Y \leq y_j\} \\
    &= P\{X = x_i\} P\{Y = y_j\} = p_{i \cdot} p_{\cdot j}.
  \end{align*}

  \((\Leftarrow)\) Suppose that there exist two sequences of non-negative numbers \(\{p_{i \cdot}\}\) and \(\{p_{\cdot j}\}\) such that \(p_{ij} = p_{i \cdot} p_{\cdot j}\) for all \(i, j\). Then for any \(x, y \in \mathbb{R}\), we have
  \[
    F(x, y) = \sum_{x_i \leq x} \sum_{y_j \leq y} p_{ij} = \sum_{x_i \leq x} \sum_{y_j \leq y} p_{i \cdot} p_{\cdot j} = \left(\sum_{x_i \leq x} p_{i \cdot}\right) \left(\sum_{y_j \leq y} p_{\cdot j}\right) = F_X(x) F_Y(y). \qedhere
  \]
\end{proof}

\begin{proposition}
  A jointly continuous bivariate random variable \((X, Y)\) with joint p.d.f.\@ \(f(x, y)\) is independent if and only if there exist two non-negative functions \(f_X(x)\) and \(f_Y(y)\) such that \(f(x, y) = f_X(x) f_Y(y)\) at any continuous point of \(f(x, y)\).
\end{proposition}
\begin{proof}
  \((\Rightarrow)\) Suppose that \((X, Y)\) is independent. Then for any continuous point \((x, y)\) of \(f(x, y)\), we have
  \[
    f(x, y) = \frac{\partial^2 F(x, y)}{\partial x \partial y} = \frac{\partial^2 F_X(x) F_Y(y)}{\partial x \partial y} = \frac{\mathrm{d} F_X(x)}{\mathrm{d}x} \cdot \frac{\mathrm{d} F_Y(y)}{\mathrm{d}y} = f_X(x) f_Y(y).
  \]

  \((\Leftarrow)\) Suppose that \(f(x, y) = f_X(x) f_Y(y)\) at any continuous point of \(f(x, y)\). Then for any \(x, y \in \mathbb{R}\), we have
  \begin{align*}
    F(x, y) &= \int_{-\infty}^x \int_{-\infty}^y f(u, v) \, \mathrm{d}v \, \mathrm{d}u = \int_{-\infty}^x \int_{-\infty}^y f_X(u) f_Y(v) \, \mathrm{d}v \, \mathrm{d}u \\
    &= \left(\int_{-\infty}^x f_X(u) \, \mathrm{d}u\right) \left(\int_{-\infty}^y f_Y(v) \, \mathrm{d}v\right) = F_X(x) F_Y(y). \qedhere
  \end{align*}
\end{proof}

\begin{proposition}
  If \((X, Y) \sim \mathrm{N}(\mu_1, \mu_2; \sigma_1^2, \sigma_2^2; \rho)\), then \(X\) and \(Y\) are independent if and only if \(\rho = 0\).
\end{proposition}

\begin{theorem}
  Let \(a, b, c\) and \(d\) be given values such that \(-\infty \leq a < b \leq \infty\) and \(-\infty \leq c < d \leq \infty\), and let \(S\) be a rectangle in the plane defined by
  \[
   S = \{(x, y) | a \leq x \leq b, c \leq y \leq d\}.
  \]
  Suppose that \(f(x, y) = 0\) for all \((x, y) \notin S\). Then the continuous (discrete) bivariate random variable \((X, Y)\) with joint p.d.f.\@ (joint p.f.\@) \(f(x, y)\) is independent if and only if there exist two non-negative functions \(f_X(x)\) and \(f_Y(y)\) such that
  \[
    f(x, y) = f_X(x) f_Y(y)
  \]
  for all \((x, y) \in S\).
\end{theorem}

\section{Functions of Bivariate Random Variables}

\begin{proposition}
  If \(X \sim \mathrm{B}(n_1, p)\) and \(Y \sim \mathrm{B}(n_2, p)\) are independent, then \(X + Y \sim \mathrm{B}(n_1 + n_2, p)\).
\end{proposition}

\begin{proposition}
  If \(X \sim \mathrm{Poi}(\lambda_1)\) and \(Y \sim \mathrm{Poi}(\lambda_2)\) are independent, then \(X + Y \sim \mathrm{Poi}(\lambda_1 + \lambda_2)\).
\end{proposition}

\paragraph{The Convolution Formula.}
Suppose that the joint p.d.f.\@ of \(X, Y\) is \(f(x, y)\). Let \(Z = X + Y\). We want to find the p.d.f.\@ of \(Z\).

The d.f.\@ of \(Z\) is given by
\[
  F_Z(z) = P(Z \leq z) = \iint_{x + y \leq z} f(x, y) \, \mathrm{d}x \, \mathrm{d}y = \int_{-\infty}^\infty \mathrm{d}x \int_{-\infty}^{z - x} f(x, y) \, \mathrm{d}y.
\]
Let \(u = y + x\), then \(\mathrm{d}u = \mathrm{d}y\). Therefore,
\[
  F_Z(z) = \int_{-\infty}^\infty \mathrm{d}x \int_{-\infty}^{z} f(x, u - x) \, \mathrm{d}u = \int_{-\infty}^z \mathrm{d}u \int_{-\infty}^\infty f(x, u - x) \, \mathrm{d}x.
\]
By differentiating \(F_Z(z)\) with respect to \(z\), we can obtain the p.d.f.\@ of \(Z\) as follows:
\[
  f_Z(z) = \int_{-\infty}^\infty f(x, z - x) \, \mathrm{d}x.
\]
By symmetry, we can also show that
\[
  f_Z(z) = \int_{-\infty}^\infty f(z - y, y) \, \mathrm{d}y.
\]

If \(X\) and \(Y\) are independent, then \(f(x, y) = f_X(x) f_Y(y)\). In this case, the p.d.f.\@ of \(Z\) can be expressed as
\[
  f_Z(z) = \int_{-\infty}^\infty f_X(x) f_Y(z - x) \, \mathrm{d}x = \int_{-\infty}^\infty f_X(z - y) f_Y(y) \, \mathrm{d}y.
\]

\begin{proposition}
  If \(X \sim \mathrm{N}(\mu_1, \sigma_1^2)\) and \(Y \sim \mathrm{N}(\mu_2, \sigma_2^2)\) are independent, then \(X + Y \sim \mathrm{N}(\mu_1 + \mu_2, \sigma_1^2 + \sigma_2^2)\).
\end{proposition}

\paragraph{The Distribution of \(M = \max\{X, Y\}\) and \(N = \min\{X, Y\}\).} Suppose that \(X\) and \(Y\) are independent with marginal d.f.\@ \(F_X(x)\) and \(F_Y(y)\) and marginal p.d.f.\@ \(f_X(x)\) and \(f_Y(y)\). We want to find the distribution of \(M = \max\{X, Y\}\) and \(N = \min\{X, Y\}\).

Since \(M = \max\{X, Y\} \leq m \Leftrightarrow X \leq m, Y \leq m\), we have
\[
  F_{M}(m) = P(M \leq m) = P(X \leq m, Y \leq m) = F_X(m) F_Y(m).
\]
Thus,
\[
  f_{M}(m) = f_X(m) F_Y(m) + f_Y(m) F_X(m).
\]

Since \(N = \min\{X, Y\} > n \Leftrightarrow X > n, Y > n\), we have
\begin{align*}
  F_{N}(n) &= P(N \leq n) = P(\min\{X, Y\} \leq n) = 1 - P(\min\{X, Y\} > n) \\
  &= 1 - P(X > n, Y > n) = 1 - [1 - F_X(n)][1 - F_Y(n)] \\
  &= F_X(n) + F_Y(n) - F_X(n) F_Y(n).
\end{align*}
Thus,
\[
  f_{N}(n) = f_X(n) [1 - F_Y(n)] + f_Y(n) [1 - F_X(n)].
\]

\section{Numerical Characteristics of Random Vectors}

\begin{theorem}
  Let \((X\) and \(Y)\) be jointly distributed random variables with p.f.\@ \(p(x, y)\) or p.d.f.\@ \(f(x, y)\). Then the expected value of a function \(h(X, Y)\), denoted by \(\mathbb{E}[h(X, Y)]\), is given by
  \[
    \mathbb{E}[h(X, Y)] = \begin{cases}
      \sum_{x} \sum_{y} h(x, y) p(x, y), & \text{if } (X, Y) \text{ is discrete}, \\
      \int_{-\infty}^\infty \int_{-\infty}^\infty h(x, y) f(x, y) \, \mathrm{d}x \, \mathrm{d}y, & \text{if } (X, Y) \text{ is continuous}.
    \end{cases}
  \]
  The expectation exists if and only if the corresponding series or integral converges absolutely.
\end{theorem}

\begin{proposition}
  \(\displaystyle \mathbb{E} \left(\sum_{i=1}^n X_i\right) = \sum_{i=1}^n \mathbb{E}(X_i)\).
\end{proposition}

\begin{proposition}
  If \(X_1, X_2, \dots, X_n\) are independent, then
  \[
    \mathbb{E} \left(\prod_{i=1}^n X_i\right) = \prod_{i=1}^n \mathbb{E}(X_i).
  \]
\end{proposition}

\begin{proposition}
  If \(X_1, X_2, \dots, X_n\) are independent, then
  \[
    \mathrm{Var} \left(\sum_{i=1}^n X_i\right) = \sum_{i=1}^n \mathrm{Var}(X_i).
  \]
\end{proposition}

\begin{definition}[Covariance]\index{covariance}
  The \emph{covariance} of two random variables \(X\) and \(Y\) is defined as
  \[
    \mathrm{Cov}(X, Y) = \mathbb{E}[(X - \mathbb{E}(X))(Y - \mathbb{E}(Y))] = \mathbb{E}(XY) - \mathbb{E}(X) \mathbb{E}(Y).
  \]
\end{definition}

\begin{proposition}
  Some properties of covariance are as follows:
  \begin{enumerate}
    \item \(\mathrm{Cov}(X, X) = \mathrm{Var}(X)\).
    \item \(\mathrm{Cov}(X, Y) = \mathrm{Cov}(Y, X)\).
    \item \(\mathrm{Cov}(X, C) = 0\) for any constant \(C\).
    \item If \(X\) and \(Y\) are independent, then \(\mathrm{Cov}(X, Y) = 0\).
    \item\label{prop:covariance:linear} \(\mathrm{Cov}(aX + b, cY + d) = ac \mathrm{Cov}(X, Y)\) for any constants \(a, b, c\) and \(d\).
    \item \(\mathrm{Cov}(X_1 + X_2, Y) = \mathrm{Cov}(X_1, Y) + \mathrm{Cov}(X_2, Y)\).
    \item \(\mathrm{Var}(X + Y) = \mathrm{Var}(X) + \mathrm{Var}(Y) + 2 \mathrm{Cov}(X, Y)\).
  \end{enumerate}
\end{proposition}

\begin{proof}[Proof of Property \ref{prop:covariance:linear}]
  We have
  \begin{align*}
    \mathrm{Cov}(aX + b, cY + d) &= \mathbb{E}[(aX + b - \mathbb{E}(aX + b))(cY + d - \mathbb{E}(cY + d))] \\
    &= \mathbb{E}[(aX - a\mathbb{E}(X))(cY - c\mathbb{E}(Y))] \\
    &= ac \mathbb{E}[(X - \mathbb{E}(X))(Y - \mathbb{E}(Y))] = ac \mathrm{Cov}(X, Y). \qedhere
  \end{align*}
\end{proof}

\begin{definition}[Correlation Coefficient]\index{correlation coefficient}
  The \emph{correlation coefficient} of two random variables \(X\) and \(Y\) is defined as
  \[
    \rho(X, Y) = \frac{\mathrm{Cov}(X, Y)}{\sqrt{\mathrm{Var}(X)} \sqrt{\mathrm{Var}(Y)}}.
  \]
\end{definition}

The correlation coefficient \(\rho(X, Y)\) is a measure of the linear relationship between \(X\) and \(Y\). It takes values in the interval \([-1, 1]\), where \(\rho(X, Y) = 1\) indicates a perfect positive linear relationship, i.e., $P\{Y = aX + b\} = 1$ for some constants \(a > 0\) and \(b\); \(\rho(X, Y) = -1\) indicates a perfect negative linear relationship, i.e., $P\{Y = aX + b\} = 1$ for some constants \(a < 0\) and \(b\); and \(\rho(X, Y) = 0\) indicates no linear relationship between \(X\) and \(Y\).

\begin{example}
  Toss a coin for 2026 times. Let \(X\) be the number of heads and \(Y\) be the number of tails. Then \(\rho(X, Y) = -1\) since \(Y = 2026 - X\).
\end{example}

\begin{example}
  Suppose that the joint p.d.f.\@ of \(X, Y\) is given by
  \[
    f(x, y) = \begin{cases}
      3x, & 0 < y < x < 1, \\
      0, & \text{otherwise}.
    \end{cases}
  \]
  Find \(\rho(X, Y)\).
\end{example}
\begin{solution}
  \begin{align*}
    \mathbb{E}(X) &= \int_{-\infty}^\infty \int_{-\infty}^\infty x f(x, y) \, \mathrm{d}x \, \mathrm{d}y = \int_0^1 3x^2 \, \mathrm{d}x \int_0^x \mathrm{d}y = \frac{3}{4}, \\
    \mathbb{E}(X^2) &= \int_{-\infty}^\infty \int_{-\infty}^\infty x^2 f(x, y) \, \mathrm{d}x \, \mathrm{d}y = \int_0^1 3x^3 \, \mathrm{d}x \int_0^x \mathrm{d}y = \frac{3}{5}, \\
    \mathbb{E}(Y) &= \int_{-\infty}^\infty \int_{-\infty}^\infty y f(x, y) \, \mathrm{d}x \, \mathrm{d}y = \int_0^1 3x \,\mathrm{d}x \int_0^x y \, \mathrm{d}y = \frac{3}{8}, \\
    \mathbb{E}(Y^2) &= \int_{-\infty}^\infty \int_{-\infty}^\infty y^2 f(x, y) \, \mathrm{d}x \, \mathrm{d}y = \int_0^1 3x \,\mathrm{d}x \int_0^x y^2 \, \mathrm{d}y = \frac{1}{5}, \\
    \mathbb{E}(XY) &= \int_{-\infty}^\infty \int_{-\infty}^\infty xy f(x, y) \, \mathrm{d}x \, \mathrm{d}y = \int_0^1 3x^2 \,\mathrm{d}x \int_0^x y \, \mathrm{d}y = \frac{3}{10}.
  \end{align*}
  Therefore,
  \begin{align*}
    \mathrm{Cov}(X, Y) &= \mathbb{E}(XY) - \mathbb{E}(X) \mathbb{E}(Y) = \frac{3}{10} - \frac{3}{4} \cdot \frac{3}{8} = \frac{3}{160}, \\
    \mathrm{Var}(X) &= \mathbb{E}(X^2) - (\mathbb{E}(X))^2 = \frac{3}{5} - \left(\frac{3}{4}\right)^2 = \frac{3}{80}, \\
    \mathrm{Var}(Y) &= \mathbb{E}(Y^2) - (\mathbb{E}(Y))^2 = \frac{1}{5} - \left(\frac{3}{8}\right)^2 = \frac{19}{320}.
  \end{align*}
  Thus,
  \[
    \rho(X, Y) = \frac{\mathrm{Cov}(X, Y)}{\sqrt{\mathrm{Var}(X)} \sqrt{\mathrm{Var}(Y)}} = \frac{\frac{3}{160}}{\sqrt{\frac{3}{80}} \sqrt{\frac{19}{320}}} = \frac{3}{\sqrt{57}}. \qedhere
  \]
\end{solution}

\begin{example}
  Suppose that \(\xi = aX + b (a > 0)\), \(\eta = cY + d (c > 0)\). Prove that \(\rho(\xi, \eta) = \rho(X, Y)\).
\end{example}
\begin{proof}
  Since \(\xi = aX + b\) and \(\eta = cY + d\), we have
  \begin{align*}
    \mathrm{Cov}(\xi, \eta) &= \mathrm{Cov}(aX + b, cY + d) = ac \mathrm{Cov}(X, Y), \\
    \mathrm{Var}(\xi) &= \mathrm{Var}(aX + b) = a^2 \mathrm{Var}(X), \\
    \mathrm{Var}(\eta) &= \mathrm{Var}(cY + d) = c^2 \mathrm{Var}(Y).
  \end{align*}
  Therefore,
  \[
    \rho(\xi, \eta) = \frac{\mathrm{Cov}(\xi, \eta)}{\sqrt{\mathrm{Var}(\xi)} \sqrt{\mathrm{Var}(\eta)}} = \frac{ac \mathrm{Cov}(X, Y)}{\sqrt{a^2 \mathrm{Var}(X)} \sqrt{c^2 \mathrm{Var}(Y)}} = \rho(X, Y). \qedhere
  \]
\end{proof}

\begin{theorem}
  If two random variables \(X\) and \(Y\) are independent, then they are uncorrelated. However, the converse is not necessarily true.
\end{theorem}

\begin{example}
  Suppose that the joint p.d.f.\@ of \(X, Y\) is given by
  \[
    f(x, y) = \begin{cases}
      \frac{1}{\pi}, & x^2 + y^2 \leq 1, \\
      0, & \text{otherwise}.
    \end{cases}
  \]
  \begin{enumerate}
    \item Prove that \(X\) and \(Y\) are uncorrelated.
    \item Prove that \(X\) and \(Y\) are not independent.
  \end{enumerate}
\end{example}
\begin{proof}
  \begin{enumerate}
    \item Evaluating the following integrals using the result of Theorem \ref{thm:important_integrals}, we have
    \begin{align*}
      \mathbb{E}(X) &= \int_{-\infty}^\infty \int_{-\infty}^\infty x f(x, y) \, \mathrm{d}x \, \mathrm{d}y = \frac{1}{\pi} \int_{-1}^1 x \, \mathrm{d}x \int_{-\sqrt{1 - x^2}}^{\sqrt{1 - x^2}} \, \mathrm{d}y = 0, \\
      \mathbb{E}(X^2) &= \int_{-\infty}^\infty \int_{-\infty}^\infty x^2 f(x, y) \, \mathrm{d}x \, \mathrm{d}y = \frac{1}{\pi} \int_{0}^{2\pi} \int_{0}^{1} r^3 \cos^2 \theta \, \mathrm{d}r \, \mathrm{d}\theta = \frac{1}{4}.
    \end{align*}
    Then,
    \[
      \mathrm{Var}(X) = \mathbb{E}(X^2) - (\mathbb{E}(X))^2 = \frac{1}{4}.
    \]
    By symmetry, we also have \(\mathbb{E}(Y) = 0\), \(\mathbb{E}(Y^2) = \frac{1}{4}\) and \(\mathrm{Var}(Y) = \frac{1}{4}\). Moreover,
    \[
      \mathbb{E}(XY) = \int_{-\infty}^\infty \int_{-\infty}^\infty xy f(x, y) \, \mathrm{d}x \, \mathrm{d}y = \frac{1}{\pi} \int_{0}^{2\pi} \int_{0}^{1} r^3 \cos \theta \sin \theta \, \mathrm{d}r \, \mathrm{d}\theta = 0.
    \]
    Therefore,
    \[
      \mathrm{Cov}(X, Y) = \mathbb{E}(XY) - \mathbb{E}(X) \mathbb{E}(Y) = 0 - 0 = 0.
    \]
    Then
    \[
      \rho(X, Y) = \frac{\mathrm{Cov}(X, Y)}{\sqrt{\mathrm{Var}(X)} \sqrt{\mathrm{Var}(Y)}} = 0,
    \]
    which means that \(X\) and \(Y\) are uncorrelated.
    \item The marginal p.d.f.\@ of \(X\) is given by
    \[
      f_X(x) = \int_{-\infty}^\infty f(x, y) \, \mathrm{d}y = \begin{cases}
        \frac{2}{\pi}\sqrt{1 - x^2}, & -1 < x < 1, \\
        0, & \text{otherwise},
      \end{cases}
    \]
    and
    \[
      f_Y(y) = \int_{-\infty}^\infty f(x, y) \, \mathrm{d}x = \begin{cases}
        \frac{2}{\pi}\sqrt{1 - y^2}, & -1 < y < 1, \\
        0, & \text{otherwise}.
      \end{cases}
    \]
    Obviously, \(f(\frac{1}{2}, \frac{1}{2}) \neq f_X(\frac{1}{2}) f_Y(\frac{1}{2})\). Thus, \(X\) and \(Y\) are not independent. \qedhere
  \end{enumerate}
\end{proof}

\begin{proposition}
  If \((X, Y) \sim \mathrm{N}(\mu_1, \mu_2; \sigma_1^2, \sigma_2^2; \rho)\), then \(\rho(X, Y) = \rho\).
\end{proposition}

\begin{definition}[Origin Moment]\index{origin moment}
  The \(n\)-th \emph{origin moment} of a random variable \(X\) is defined as
  \[
    m_n = \mathbb{E}(X^n).
  \]
  The \(n\)-th \emph{mixed origin moment} of two random variables \(X\) and \(Y\) is defined as
  \[
    m_{n_1, n_2} = \mathbb{E}(X^{n_1} Y^{n_2}).
  \]
\end{definition}

\begin{definition}[Central Moment]\index{central moment}
  The \(n\)-th \emph{central moment} of a random variable \(X\) is defined as
  \[
    \mu_n = \mathbb{E}[(X - \mathbb{E}(X))^n].
  \]
  The \(n\)-th \emph{mixed central moment} of two random variables \(X\) and \(Y\) is defined as
  \[
    \mu_{n_1, n_2} = \mathbb{E}[(X - \mathbb{E}(X))^{n_1} (Y - \mathbb{E}(Y))^{n_2}].
  \]
\end{definition}

\begin{definition}[Covariance Matrix]\index{covariance matrix}
  For a random vector \(\mathbf{X} = (X_1, X_2, \dots, X_n)\), the \emph{covariance matrix} of \(\mathbf{X}\) is defined as
  \[
    \Sigma = \begin{bmatrix}
      c_{11} & c_{12} & \cdots & c_{1n} \\
      c_{21} & c_{22} & \cdots & c_{2n} \\
      \vdots & \vdots & \ddots & \vdots \\
      c_{n1} & c_{n2} & \cdots & c_{nn}
    \end{bmatrix},
  \]
  if all the covariances \(c_{ij} = \mathrm{Cov}(X_i, X_j)\) exist for \(i, j = 1, 2, \dots, n\).
\end{definition}

\begin{definition}[Multivariate Normal Distribution]\index{multivariate normal distribution}
  A random vector \(\mathbf{X} = (X_1, X_2, \dots, X_n)\) with mean vector \(\mu\) and covariance matrix \(\Sigma\) is said to have a \emph{multivariate normal distribution} if its joint p.d.f.\@ is given by
  \[
    f_{\mathbf{X}}(\mathbf{x}) = \frac{1}{(2\pi)^{\frac{n}{2}} |\Sigma|^{\frac{1}{2}}} \exp\left(-\frac{1}{2} (\mathbf{x} - \mu)^T \Sigma^{-1} (\mathbf{x} - \mu)\right),
  \]
  where \(\mathbf{x} = (x_1, x_2, \dots, x_n)\) and \(|\Sigma|\) is the determinant of \(\Sigma\).  
\end{definition}

\begin{proposition}\label{prop:multivariate-normal-linear-combination}
  Suppose that \(X_1, X_2, \dots, X_n\) are \(n\) random variables. Then \((X_1, X_2, \dots, X_n)\) has a multivariate normal distribution if and only if any linear combination of \(X_1, X_2, \dots, X_n\) has a normal distribution.
\end{proposition}

\begin{proposition}\label{prop:multivariate-normal-linear-transformation}
  If \(\mathbf{Y} = \mathbf{A}\mathbf{X} + b\), where \(\mathbf{A}\) is an \(n \times n\) non-singular matrix and \(b\) is a (column) \(n\)-dimensional vector of constants, then \(\mathbf{Y} \sim \mathrm{N}(\mathbf{A}\mu + b, \mathbf{A} \Sigma \mathbf{A}^T)\), where \(\mu = \mathbb{E}(\mathbf{X})\) and \(\Sigma\) is the covariance matrix of the multivariate normal random variable \(\mathbf{X}\).
\end{proposition}

\begin{proposition}
  If \(\mathbf{X} = (X_1, X_2, \dots, X_n)\) has \(n\) dimensional normal distribution, then \(X_j (j = 1, 2, \dots, n)\) has a normal distribution. Conversely, if \(X_1, X_2, \dots, X_n\) have normal distributions and are independent, then \((X_1, X_2, \dots, X_n)\) has \(n\) dimensional multivariate normal distribution.
\end{proposition}

\begin{proposition}
  Suppose that \((X_1, X_2, \dots, X_n)\) has a multivariate normal distribution. Then \(X_1, X_2, \dots, X_n\) are independent if and only if \(X_1, X_2, \dots, X_n\) are pairwise uncorrelated.
\end{proposition}

\begin{example}
  Suppose that \(X \sim \mathrm{N}(1, 2)\), \(Y \sim \mathrm{N}(0, 1)\) and \(X\) and \(Y\) are independent. Find the p.d.f.\@ of \(Z = 2X - Y + 3\).
\end{example}
\begin{solution}
  Since \(X \sim \mathrm{N}(1, 2)\), \(Y \sim \mathrm{N}(0, 1)\) and \(X\) and \(Y\) are independent, we have
  \begin{align*}
    \mathbb{E}(Z) &= 2 \mathbb{E}(X) - \mathbb{E}(Y) + 3 = 2 \cdot 1 - 0 + 3 = 5, \\
    \mathrm{Var}(Z) &= 2^2 \mathrm{Var}(X) + (-1)^2 \mathrm{Var}(Y) = 2^2 \cdot 2 + (-1)^2 \cdot 1 = 3^3.
  \end{align*}
  By Proposition \ref{prop:multivariate-normal-linear-combination}, \(Z\) has a normal distribution, i.e.,
  \[
    Z = 2X - Y + 3 \sim \mathrm{N}(5, 3^3).
  \]
  Thus, the p.d.f.\@ of \(Z\) is given by
  \[
    f_Z(z) = \frac{1}{3\sqrt{2 \pi}} e^{-\frac{(z - 5)^2}{18}}. \qedhere
  \]
\end{solution}

\begin{example}
  Suppose that \(X \sim \mathrm{N}(0, 1)\), \(Y \sim \mathrm{N}(0, 1)\) and \(X\) and \(Y\) are independent. Calculate the value of \(\mathbb{E}(|X - Y|)\).
\end{example}
\begin{solution}
  Let \(Z = X - Y\). Then \(Z \sim \mathrm{N}(0, 2)\), i.e.,
  \[
    f_Z(z) = \frac{1}{2\sqrt{\pi}} e^{-\frac{z^2}{4}}.
  \]
  Therefore,
  \[
    \mathbb{E}(|X - Y|) = \mathbb{E}(|Z|) = \int_{-\infty}^\infty |z| \frac{1}{2\sqrt{\pi}} e^{-\frac{z^2}{4}} \, \mathrm{d}z = \frac{1}{\sqrt{\pi}} \int_0^\infty z e^{-\frac{z^2}{4}} \, \mathrm{d}z = \frac{2}{\sqrt{\pi}}. \qedhere
  \]
\end{solution}

\begin{example}
  Suppose that \((X, Y) \sim \mathrm{N}(0, 1, 3^2, 4^2, -\frac{1}{2})\), \(Z = \frac{X}{3} + \frac{Y}{2}\). Prove that \(X\) and \(Z\) are independent.
\end{example}
\begin{solution}
  Since \((X, Y) \sim \mathrm{N}(0, 1, 3^2, 4^2, -\frac{1}{2})\), we have \(\mathrm{Var}(X) = 3^2\), \(\mathrm{Var}(Y) = 4^2\) and \(\rho(X, Y) = -\frac{1}{2}\). Note that
  \[
    \begin{bmatrix}
      X \\
      Z
    \end{bmatrix} = \begin{bmatrix}
      1 & 0 \\
      \frac{1}{3} & \frac{1}{2}
    \end{bmatrix} \begin{bmatrix}
      X \\
      Y
    \end{bmatrix}.
  \]
  By Proposition \ref{prop:multivariate-normal-linear-transformation}, \((X, Z)\) has a bivariate normal distribution. Moreover,
  \begin{align*}
    \mathrm{Cov}(X, Z) &= \mathrm{Cov}\!\left(X, \frac{X}{3} + \frac{Y}{2}\right) = \frac{1}{3} \mathrm{Cov}(X) + \frac{1}{2} \mathrm{Cov}(X, Y) \\
    &= \frac{1}{3} \cdot 3^2 + \frac{1}{2} \cdot (-\frac{1}{2}) \cdot 3 \cdot 4 = 0.
  \end{align*}
  Thus, \(X\) and \(Z\) are uncorrelated. Since \((X, Z)\) has a bivariate normal distribution, \(X\) and \(Z\) are independent. \qedhere
\end{solution}